\section{Fachanwender}

Handbuch f\textbackslash\{\}"ur die t\textbackslash\{\}"agliche Nutzung der Desktop-App ohne technische Implementierungsdetails. Begriffe und Schritte sind an der realen Oberfl\textbackslash\{\}"ache und den Modulen unter \texttt{docs\_manual/modules/} ausgerichtet.

\subsection{Aufgaben}

\begin{itemize}
\item Im Bereich \textbf{Operations} Chats f\textbackslash\{\}"uhren, Inhalte mit dem Modell erarbeiten.
\item Unter \textbf{Knowledge} Dokumente verwalten und f\textbackslash\{\}"ur RAG-basierte Antworten nutzen (wenn aktiviert).
\item Im \textbf{Prompt Studio} Vorlagen pflegen und im Arbeitsfluss verwenden.
\item \textbf{Agent Tasks} und \textbf{Agentenwahl im Chat} nutzen, wenn Ihre Umgebung das vorsieht.
\item \textbf{Einstellungen} \textbackslash\{\}"offnen: Erscheinungsbild, KI/Modelle, ggf. erweiterte Optionen inkl. \textbf{Chat-Kontext-Modus}.
\item \textbackslash\{\}"Uber \textbf{Control Center} pr\textbackslash\{\}"ufen, ob Modelle und Provider erreichbar sind, bevor Sie chatten.
\end{itemize}

\subsection{Typische Workflows}

\subsubsection{Neuen Chat starten und schreiben}

\begin{enumerate}
\item Sidebar: \textbf{Operations} w\textbackslash\{\}"ahlen.
\item In der linken Leiste \textbf{Chat} aktivieren (Workspace \textbf{Chat}).
\item Session w\textbackslash\{\}"ahlen oder neue Session anlegen.
\item Unten Text eingeben und senden.
\item Modell oder Agent \textbackslash\{\}"uber die Kopfzeile des Chats w\textbackslash\{\}"ahlen, soweit angeboten.
\end{enumerate}

\subsubsection{Rolle oder Routing per Eingabezeile}

Am Anfang der Nachricht, direkt nach den Regeln der App:

\begin{itemize}
\item \texttt{/think \textbackslash\{\}ldots\{\}}, \texttt{/code \textbackslash\{\}ldots\{\}}, \texttt{/research \textbackslash\{\}ldots\{\}} usw. setzen die \textbf{Rolle} f\textbackslash\{\}"ur diese Nachricht (Rest der Zeile wird mitgesendet).
\item \texttt{/auto on} oder \texttt{/auto off} schalten \textbf{Auto-Routing}.
\item \texttt{/cloud on} oder \texttt{/cloud off} steuern \textbf{Cloud-Eskalation} (nur wirksam, wenn Ihre Organisation Ollama-Cloud und Schl\textbackslash\{\}"ussel freigibt).
\item \texttt{/delegate Aufgabe} startet die \textbf{Delegations}-Bearbeitung mit dem Text nach dem Befehl --- ohne Text erscheint nur ein Hinweis.
\end{itemize}

(Implementierung: \texttt{app/core/commands/chat\_commands.py}.)

\subsubsection{RAG nutzen}

\begin{enumerate}
\item \textbf{Operations \$\textbackslash\{\}rightarrow\$ Knowledge:} Quellen anlegen oder aktualisieren, Indexierung ausf\textbackslash\{\}"uhren (wie in Ihrer Umgebung vorgesehen, z. B. \textbackslash\{\}"uber Skripte aus dem Projekt).
\item In den \textbf{Einstellungen} RAG aktivieren und \textbf{Space} sowie \textbf{Top-K} pr\textbackslash\{\}"ufen (\texttt{rag\_enabled}, \texttt{rag\_space}, \texttt{rag\_top\_k}).
\item Im \textbf{Chat} Frage stellen, die zu den indexierten Inhalten passt.
\end{enumerate}

\subsubsection{Kontext verstehen (ohne Technik)}

Der \textbf{Chat-Kontext} ist zus\textbackslash\{\}"atzlicher Hintergrund, den die App dem Modell mitgeben kann: z. B. zu welchem \textbf{Projekt} und welchem \textbf{Chat} die Nachricht geh\textbackslash\{\}"ort. Dar\textbackslash\{\}"uber gibt die Leiste \textbf{\textbackslash\{\}"uber} der Konversation Auskunft (Projektname, Chat-Titel, ggf. Topic).

\textbf{Drei Stufen ?wie viel Hintergrund`` (vereinfacht):}

\begin{itemize}
\item \textbf{Aus:} Es wird kein solcher Struktur-Hintergrund mitgeschickt. Das Modell sieht nur das, was Sie in den Nachrichten schreiben (plus ggf. RAG, Prompts, Agenten-Systemprompt).
\item \textbf{Neutral:} Hintergrund wird sachlich/knapp formuliert mitgegeben.
\item \textbf{Semantisch:} Hintergrund wird ausf\textbackslash\{\}"uhrlicher mitgegeben.
\end{itemize}

Wo Sie das umstellen: \textbf{Einstellungen \$\textbackslash\{\}rightarrow\$ Advanced \$\textbackslash\{\}rightarrow\$ Chat-Kontext-Modus} (\texttt{AdvancedSettingsPanel}).

\textbf{Wichtig:} Wenn Antworten ?ohne Bezug zum Projekt\texttt{} wirken, kann der Modus \textbf{Aus} aktiv sein --- oder eine andere Regel in der App hat Vorrang vor Ihrer manuellen Einstellung (das kl\textbackslash\{\}"aren ggf. Admins oder Support \textbackslash\{\}"uber die QA-/Debug-Ansichten).

\subsubsection{Typische Fehler und was Sie tun k\textbackslash\{\}"onnen}

\begin{center}
\begin{tabular}{ll}
\hline
Symptom & Was pr\textbackslash\{\}"ufen \\
\hline
Keine Antwort / dauert sehr lange & \textbf{Control Center \$\textbackslash\{\}rightarrow\$ Providers:} Online? Ollama auf dem Rechner/Server gestartet? \\
Modellliste leer & Modell mit \texttt{ollama pull \textbackslash\{\}ldots\{\}} laden; App neu starten; \textbf{Einstellungen \$\textbackslash\{\}rightarrow\$ AI / Models} oder \textbf{Control Center \$\textbackslash\{\}rightarrow\$ Models}. \\
RAG liefert nichts & Unter \textbf{Knowledge} indexiert? RAG in den Einstellungen eingeschaltet? Richtiger \textbf{Space}? \\
Slash-Befehl wird als normaler Text gesendet & Zeile muss mit \texttt{/} beginnen; Tippfehler im Befehlsnamen \$\textbackslash\{\}rightarrow\$ App meldet ?Unbekannter Befehl\texttt{}. \\
\texttt{/delegate} passiert nichts Sichtbares & Kein Text nach \texttt{/delegate} eingegeben --- es erscheint nur die Verwendungsmeldung. \\
Chat wirkt ?ohne Projektbezug\texttt{} & Kontextmodus \textbf{Aus}; oder Projekt/Chat in der Kontextleiste pr\textbackslash\{\}"ufen. \\
\hline
\end{tabular}
\end{center}

\subsection{Genutzte Module}

\begin{center}
\begin{tabular}{ll}
\hline
Modul & Nutzen f\textbackslash\{\}"ur Sie \\
\hline
gui (\texttt{../../modules/gui/README.md}) & Navigation, Operations, Chat-, Knowledge-, Prompt-Studio-Workspaces. \\
chat (\texttt{../../modules/chat/README.md}) & Senden, Streaming, Slash-Commands. \\
context (\texttt{../../modules/context/README.md}) & Bedeutung von Kontextmodus (fachlich vereinfacht oben). \\
prompts (\texttt{../../modules/prompts/README.md}) & Prompt Studio. \\
rag (\texttt{../../modules/rag/README.md}) & Knowledge-Workspace und RAG-Schalter. \\
agents (\texttt{../../modules/agents/README.md}) & Agent Tasks, Agentenwahl, \texttt{/delegate}. \\
settings (\texttt{../../modules/settings/README.md}) & Alle Einstellungskategorien inkl. Advanced. \\
providers (\texttt{../../modules/providers/README.md}) & Providers-Workspace (Status). \\
\hline
\end{tabular}
\end{center}

\subsection{Risiken}

\begin{itemize}
\item \textbf{Falsche Annahmen zum Kontext:} Bei Modus \textbf{Aus} oder eingeschr\textbackslash\{\}"ankten Includes ?wei\textbackslash\{\}ss\{\}\texttt{} das Modell nicht automatisch alles \textbackslash\{\}"uber Projekt/Chat.
\item \textbf{Cloud und Daten:} Wenn \textbf{Cloud} eingeschaltet ist und Ihre Organisation Cloud-Modelle nutzt, verlassen Anfragen das lokale Ollama-Szenario --- Kl\textbackslash\{\}"arung mit Admin/Richtlinien.
\item \textbf{RAG-Qualit\textbackslash\{\}"at:} Antworten h\textbackslash\{\}"angen von Index, gew\textbackslash\{\}"ahltem Space und Dokumentqualit\textbackslash\{\}"at ab.
\end{itemize}

\subsection{Best Practices}

\begin{itemize}
\item Vor l\textbackslash\{\}"angeren Sessions \textbf{Provider-Status} und \textbf{Standardmodell} pr\textbackslash\{\}"ufen.
\item F\textbackslash\{\}"ur wiederkehrende Texte \textbf{Prompt Studio} nutzen statt jedes Mal neu zu formulieren.
\item Bei Themen mit internen Dokumenten \textbf{Knowledge} pflegen und RAG gezielt aktivieren.
\item Wenn Antworten zu allgemein sind: Kontextmodus nicht auf \textbf{Aus}, sofern Ihre Richtlinien Projektbezug erlauben.
\end{itemize}
